\IfFileExists{stacks-project.cls}{%
\documentclass{stacks-project}
}{%
\documentclass{amsart}
}

% For dealing with references we use the comment environment
\usepackage{verbatim}
\newenvironment{reference}{\comment}{\endcomment}
%\newenvironment{reference}{}{}
\newenvironment{slogan}{\comment}{\endcomment}
\newenvironment{history}{\comment}{\endcomment}

% For commutative diagrams we use Xy-pic
\usepackage[all]{xy}

% We use 2cell for 2-commutative diagrams.
\xyoption{2cell}
\UseAllTwocells

% We use multicol for the list of chapters between chapters
\usepackage{multicol}

% This is generally recommended for better output
\usepackage{lmodern}
\usepackage[T1]{fontenc}

% For cross-file-references
\usepackage{xr-hyper}

% Package for hypertext links:
\usepackage{hyperref}

% For any local file, say "hello.tex" you want to link to please
% use \externaldocument[hello-]{hello}
\externaldocument[introduction-]{introduction}
\externaldocument[conventions-]{conventions}
\externaldocument[sets-]{sets}
\externaldocument[categories-]{categories}
\externaldocument[topology-]{topology}
\externaldocument[sheaves-]{sheaves}
\externaldocument[sites-]{sites}
\externaldocument[stacks-]{stacks}
\externaldocument[fields-]{fields}
\externaldocument[algebra-]{algebra}
\externaldocument[brauer-]{brauer}
\externaldocument[homology-]{homology}
\externaldocument[derived-]{derived}
\externaldocument[simplicial-]{simplicial}
\externaldocument[more-algebra-]{more-algebra}
\externaldocument[smoothing-]{smoothing}
\externaldocument[modules-]{modules}
\externaldocument[sites-modules-]{sites-modules}
\externaldocument[injectives-]{injectives}
\externaldocument[cohomology-]{cohomology}
\externaldocument[sites-cohomology-]{sites-cohomology}
\externaldocument[dga-]{dga}
\externaldocument[dpa-]{dpa}
\externaldocument[sdga-]{sdga}
\externaldocument[hypercovering-]{hypercovering}
\externaldocument[schemes-]{schemes}
\externaldocument[constructions-]{constructions}
\externaldocument[properties-]{properties}
\externaldocument[morphisms-]{morphisms}
\externaldocument[coherent-]{coherent}
\externaldocument[divisors-]{divisors}
\externaldocument[limits-]{limits}
\externaldocument[varieties-]{varieties}
\externaldocument[topologies-]{topologies}
\externaldocument[descent-]{descent}
\externaldocument[perfect-]{perfect}
\externaldocument[more-morphisms-]{more-morphisms}
\externaldocument[flat-]{flat}
\externaldocument[groupoids-]{groupoids}
\externaldocument[more-groupoids-]{more-groupoids}
\externaldocument[etale-]{etale}
\externaldocument[chow-]{chow}
\externaldocument[intersection-]{intersection}
\externaldocument[pic-]{pic}
\externaldocument[weil-]{weil}
\externaldocument[adequate-]{adequate}
\externaldocument[dualizing-]{dualizing}
\externaldocument[duality-]{duality}
\externaldocument[discriminant-]{discriminant}
\externaldocument[derham-]{derham}
\externaldocument[local-cohomology-]{local-cohomology}
\externaldocument[algebraization-]{algebraization}
\externaldocument[curves-]{curves}
\externaldocument[resolve-]{resolve}
\externaldocument[models-]{models}
\externaldocument[functors-]{functors}
\externaldocument[equiv-]{equiv}
\externaldocument[pione-]{pione}
\externaldocument[etale-cohomology-]{etale-cohomology}
\externaldocument[proetale-]{proetale}
\externaldocument[relative-cycles-]{relative-cycles}
\externaldocument[more-etale-]{more-etale}
\externaldocument[trace-]{trace}
\externaldocument[crystalline-]{crystalline}
\externaldocument[spaces-]{spaces}
\externaldocument[spaces-properties-]{spaces-properties}
\externaldocument[spaces-morphisms-]{spaces-morphisms}
\externaldocument[decent-spaces-]{decent-spaces}
\externaldocument[spaces-cohomology-]{spaces-cohomology}
\externaldocument[spaces-limits-]{spaces-limits}
\externaldocument[spaces-divisors-]{spaces-divisors}
\externaldocument[spaces-over-fields-]{spaces-over-fields}
\externaldocument[spaces-topologies-]{spaces-topologies}
\externaldocument[spaces-descent-]{spaces-descent}
\externaldocument[spaces-perfect-]{spaces-perfect}
\externaldocument[spaces-more-morphisms-]{spaces-more-morphisms}
\externaldocument[spaces-flat-]{spaces-flat}
\externaldocument[spaces-groupoids-]{spaces-groupoids}
\externaldocument[spaces-more-groupoids-]{spaces-more-groupoids}
\externaldocument[bootstrap-]{bootstrap}
\externaldocument[spaces-pushouts-]{spaces-pushouts}
\externaldocument[spaces-chow-]{spaces-chow}
\externaldocument[groupoids-quotients-]{groupoids-quotients}
\externaldocument[spaces-more-cohomology-]{spaces-more-cohomology}
\externaldocument[spaces-simplicial-]{spaces-simplicial}
\externaldocument[spaces-duality-]{spaces-duality}
\externaldocument[formal-spaces-]{formal-spaces}
\externaldocument[restricted-]{restricted}
\externaldocument[spaces-resolve-]{spaces-resolve}
\externaldocument[formal-defos-]{formal-defos}
\externaldocument[defos-]{defos}
\externaldocument[cotangent-]{cotangent}
\externaldocument[examples-defos-]{examples-defos}
\externaldocument[algebraic-]{algebraic}
\externaldocument[examples-stacks-]{examples-stacks}
\externaldocument[stacks-sheaves-]{stacks-sheaves}
\externaldocument[criteria-]{criteria}
\externaldocument[artin-]{artin}
\externaldocument[quot-]{quot}
\externaldocument[stacks-properties-]{stacks-properties}
\externaldocument[stacks-morphisms-]{stacks-morphisms}
\externaldocument[stacks-limits-]{stacks-limits}
\externaldocument[stacks-cohomology-]{stacks-cohomology}
\externaldocument[stacks-perfect-]{stacks-perfect}
\externaldocument[stacks-introduction-]{stacks-introduction}
\externaldocument[stacks-more-morphisms-]{stacks-more-morphisms}
\externaldocument[stacks-geometry-]{stacks-geometry}
\externaldocument[moduli-]{moduli}
\externaldocument[moduli-curves-]{moduli-curves}
\externaldocument[examples-]{examples}
\externaldocument[exercises-]{exercises}
\externaldocument[guide-]{guide}
\externaldocument[desirables-]{desirables}
\externaldocument[coding-]{coding}
\externaldocument[obsolete-]{obsolete}
\externaldocument[fdl-]{fdl}
\externaldocument[index-]{index}

% Theorem environments.
%
\theoremstyle{plain}
\newtheorem{theorem}[subsection]{Theorem}
\newtheorem{proposition}[subsection]{Proposition}
\newtheorem{lemma}[subsection]{Lemma}

\theoremstyle{definition}
\newtheorem{definition}[subsection]{Definition}
\newtheorem{example}[subsection]{Example}
\newtheorem{exercise}[subsection]{Exercise}
\newtheorem{situation}[subsection]{Situation}

\theoremstyle{remark}
\newtheorem{remark}[subsection]{Remark}
\newtheorem{remarks}[subsection]{Remarks}

\numberwithin{equation}{subsection}

% Macros
%
\def\lim{\mathop{\mathrm{lim}}\nolimits}
\def\colim{\mathop{\mathrm{colim}}\nolimits}
\def\Spec{\mathop{\mathrm{Spec}}}
\def\Hom{\mathop{\mathrm{Hom}}\nolimits}
\def\Ext{\mathop{\mathrm{Ext}}\nolimits}
\def\SheafHom{\mathop{\mathcal{H}\!\mathit{om}}\nolimits}
\def\SheafExt{\mathop{\mathcal{E}\!\mathit{xt}}\nolimits}
\def\Sch{\mathit{Sch}}
\def\Mor{\mathop{\mathrm{Mor}}\nolimits}
\def\Ob{\mathop{\mathrm{Ob}}\nolimits}
\def\Sh{\mathop{\mathit{Sh}}\nolimits}
\def\NL{\mathop{N\!L}\nolimits}
\def\CH{\mathop{\mathrm{CH}}\nolimits}
\def\proetale{{pro\text{-}\acute{e}tale}}
\def\etale{{\acute{e}tale}}
\def\QCoh{\mathit{QCoh}}
\def\Ker{\mathop{\mathrm{Ker}}}
\def\Im{\mathop{\mathrm{Im}}}
\def\Coker{\mathop{\mathrm{Coker}}}
\def\Coim{\mathop{\mathrm{Coim}}}

% Boxtimes
%
\DeclareMathSymbol{\boxtimes}{\mathbin}{AMSa}{"02}

%
% Macros for moduli stacks/spaces
%
\def\QCohstack{\mathcal{QC}\!\mathit{oh}}
\def\Cohstack{\mathcal{C}\!\mathit{oh}}
\def\Spacesstack{\mathcal{S}\!\mathit{paces}}
\def\Quotfunctor{\mathrm{Quot}}
\def\Hilbfunctor{\mathrm{Hilb}}
\def\Curvesstack{\mathcal{C}\!\mathit{urves}}
\def\Polarizedstack{\mathcal{P}\!\mathit{olarized}}
\def\Complexesstack{\mathcal{C}\!\mathit{omplexes}}
% \Pic is the operator that assigns to X its picard group, usage \Pic(X)
% \Picardstack_{X/B} denotes the Picard stack of X over B
% \Picardfunctor_{X/B} denotes the Picard functor of X over B
\def\Pic{\mathop{\mathrm{Pic}}\nolimits}
\def\Picardstack{\mathcal{P}\!\mathit{ic}}
\def\Picardfunctor{\mathrm{Pic}}
\def\Deformationcategory{\mathcal{D}\!\mathit{ef}}
\title{The localization category of a left multiplicative system is well defined}
\author{The Stacks Project Authors}
\date{stacks-project:categories.tex:local-snapshot-20261012; GNU FDL 1.2 or later}
\begin{document}
\maketitle
\noindent\textit{Editorial note.} Counted span is the author's own proof environment, source lines 4231-4428, following the statement at 4218-4229.\par\medskip
\noindent\textit{Editorial note.} Author omits the mirror half p1 q = p2 q of part (2) (marked 'It is similar to the case we did') and leaves the identity axioms of part (3) to the reader; both are recorded here, and the executed material covers every structural step.\par\medskip
\noindent\textit{Editorial note.} The footnote-free main line uses only LMS1-LMS3 and the in-span definitions.\par\medskip
\noindent\textit{Editorial note.} Steps the author leaves to the reader are flagged here for the reviewer: source line(s) 4402 carry an "omitted"/"follows easily"/"immediate from the definitions" phrase. Those sentences are retained verbatim; the displayed argument around them stays inside the counted span.\par\medskip
\noindent\textit{Allowed standard prerequisites (editorial).}\begin{itemize}
\item Axioms LMS1-LMS3 of a left multiplicative system. Definition of a left multiplicative system in this chapter, stated in the left calculus of fractions section immediately before the lemma. Supplies the framework for all three parts: LMS2 produces the common lower-vertex diagrams, LMS3 the post-composition making two parallel arrows equal, LMS1 keeps composites of S-arrows in S.
\item Pairs, the relation E and the composition rule. Given in the text immediately preceding the lemma (construction of the objects and morphisms of S\textasciicircum{}\{-1\}C). The objects whose equivalence relation, well-definedness and associativity are verified; the proof works directly with these in-span definitions.
\end{itemize}\medskip
\begin{lemma}
\label{lemma-left-localization}
Let $\mathcal{C}$ be a category and let $S$ be a left multiplicative
system.
\begin{enumerate}
\item The relation on pairs defined above is an equivalence relation.
\item The composition rule given above is well defined on equivalence
classes.
\item Composition is associative (and the identity morphisms satisfy
the identity axioms), and hence $S^{-1}\mathcal{C}$ is a category.
\end{enumerate}
\end{lemma}

\begin{proof}
Proof of (1). Let us say two pairs $p_1 = (f_1 : X \to Y_1, s_1 : Y \to Y_1)$
and $p_2 = (f_2 : X \to Y_2, s_2 : Y \to Y_2)$ are elementary equivalent
if there exists a morphism $a : Y_1 \to Y_2$ of $\mathcal{C}$ such that
$a \circ f_1 = f_2$ and $a \circ s_1 = s_2$. Diagram:
$$
\xymatrix{
X \ar@{=}[d] \ar[r]_{f_1} & Y_1 \ar[d]^a & Y \ar[l]^{s_1} \ar@{=}[d] \\
X \ar[r]^{f_2} & Y_2 & Y \ar[l]_{s_2}
}
$$
Let us denote this property by saying $p_1Ep_2$.
Note that $pEp$ and $aEb, bEc \Rightarrow aEc$.
(Despite its name, $E$ is not an equivalence
relation.)
Part (1) claims that the relation
$p \sim p' \Leftrightarrow \exists q: pEq \wedge p'Eq$
(where $q$ is supposed to be a pair satisfying the
same conditions as $p$ and $p'$)
is an equivalence relation. A simple formal argument, using the properties
of $E$ above, shows that it suffices to prove
$p_3Ep_1, p_3Ep_2 \Rightarrow p_1 \sim p_2$.
Thus suppose that we are given a commutative diagram
$$
\xymatrix{
 & Y_1 & \\
X \ar[ru]^{f_1} \ar[r]^{f_3} \ar[rd]_{f_2} &
Y_3 \ar[u]_{a_{31}} \ar[d]^{a_{32}} &
Y \ar[lu]_{s_1} \ar[l]_{s_3} \ar[ld]^{s_2} \\
& Y_2 &
}
$$
with $s_i \in S$.
First we apply LMS2 to get a commutative diagram
$$
\xymatrix{
Y \ar[d]_{s_1} \ar[r]_{s_2} & Y_2 \ar@{..>}[d]^{a_{24}} \\
Y_1 \ar@{..>}[r]^{a_{14}} & Y_4
}
$$
with $a_{24} \in S$. Then, we have
$$
a_{14} \circ a_{31} \circ s_3 =
a_{14} \circ s_1 =
a_{24} \circ s_2 =
a_{24} \circ a_{32} \circ s_3.
$$
Hence, by LMS3, there exists a
morphism $s_{44} : Y_4 \to Y'_4$ such that $s_{44} \in S$ and
$s_{44} \circ a_{14} \circ a_{31}
= s_{44} \circ a_{24} \circ a_{32}$.
Hence, after replacing $Y_4$, $a_{14}$ and $a_{24}$ by $Y'_4$,
$s_{44} \circ a_{14}$ and $s_{44} \circ a_{24}$, we may assume
that $a_{14} \circ a_{31} = a_{24} \circ a_{32}$ (and
we still have $a_{24} \in S$ and
$a_{14} \circ s_1 = a_{24} \circ s_2$). Set
$$
f_4 =
a_{14} \circ f_1 =
a_{14} \circ a_{31} \circ f_3 =
a_{24} \circ a_{32} \circ f_3 =
a_{24} \circ f_2
$$
and
$s_4 = a_{14} \circ s_1 = a_{24} \circ s_2$. Then, the diagram
$$
\xymatrix{
X \ar@{=}[d] \ar[r]_{f_1} & Y_1 \ar[d]^{a_{14}} & Y \ar[l]^{s_1} \ar@{=}[d] \\
X \ar[r]^{f_4} & Y_4 & Y \ar[l]_{s_4}
}
$$
commutes, and we have $s_4 \in S$ (by LMS1). Thus, $p_1 E p_4$,
where $p_4 = (f_4, s_4)$. Similarly, $p_2 E p_4$. Combining these,
we find $p_1 \sim p_2$.

\medskip\noindent
Proof of (2). Let $p = (f : X \to Y', s : Y \to Y')$ and
$q = (g : Y \to Z', t : Z \to Z')$ be pairs as in the definition of composition
above. To compose we choose a diagram
$$
\xymatrix{
Y \ar[d]_s \ar[r]_g & Z' \ar[d]^{u_2} \\
Y' \ar[r]^{h_2} & Z_2
}
$$
with $u_2 \in S$. We first show that the equivalence class of the pair
$r_2 = (h_2 \circ f : X \to Z_2, u_2 \circ t : Z \to Z_2)$
is independent of the choice of $(Z_2, h_2, u_2)$. Namely, suppose
that $(Z_3, h_3, u_3)$ is another choice with corresponding composition
$r_3 = (h_3 \circ f : X \to Z_3, u_3 \circ t : Z \to Z_3)$.
Then by LMS2 we can choose a diagram
$$
\xymatrix{
Z' \ar[d]_{u_2} \ar[r]_{u_3} & Z_3 \ar[d]^{u_{34}} \\
Z_2 \ar[r]^{h_{24}} & Z_4
}
$$
with $u_{34} \in S$. We have $h_2 \circ s = u_2 \circ g$
and similarly $h_3 \circ s = u_3 \circ g$. Now,
$$
u_{34} \circ h_3 \circ s
= u_{34} \circ u_3 \circ g
= h_{24} \circ u_2 \circ g
= h_{24} \circ h_2 \circ s.
$$
Hence, LMS3 shows that there
exist a $Z'_4$ and an $s_{44} : Z_4 \to Z'_4$ such that
$s_{44} \circ u_{34} \circ h_3 = s_{44} \circ h_{24} \circ h_2$.
Replacing $Z_4$, $h_{24}$ and $u_{34}$ by $Z'_4$,
$s_{44} \circ h_{24}$ and $s_{44} \circ u_{34}$, we may
assume that $u_{34} \circ h_3 = h_{24} \circ h_2$.
Meanwhile, the relations $u_{34} \circ u_3 = h_{24} \circ u_2$
and $u_{34} \in S$ continue to hold. We can now set
$h_4 = u_{34} \circ h_3 = h_{24} \circ h_2$ and
$u_4 = u_{34} \circ u_3 = h_{24} \circ u_2$. Then, we have a
commutative diagram
$$
\xymatrix{
X \ar@{=}[d] \ar[r]_{h_2\circ f} &
Z_2 \ar[d]^{h_{24}} &
Z \ar[l]^{u_2 \circ t} \ar@{=}[d] \\
X \ar@{=}[d] \ar[r]^{h_4\circ f} &
Z_4 &
Z \ar@{=}[d] \ar[l]_{u_4 \circ t} \\
X \ar[r]^{h_3 \circ f} &
Z_3 \ar[u]^{u_{34}} &
Z \ar[l]_{u_3 \circ t}
}
$$
Hence we obtain a pair
$r_4 =
(h_4 \circ f : X \to Z_4, u_4 \circ t : Z \to Z_4)$
and the above diagram shows that we have $r_2Er_4$ and
$r_3Er_4$, whence $r_2 \sim r_3$, as desired. Thus it now
makes sense to define $p \circ q$ as the equivalence class of
all possible pairs $r$ obtained as above.

\medskip\noindent
To finish the proof of (2) we have to show that given pairs
$p_1, p_2, q$ such that $p_1Ep_2$ then $p_1 \circ q = p_2 \circ q$ and
$q \circ p_1 = q \circ p_2$ whenever the compositions make sense.
To do this, write $p_1 = (f_1 : X \to Y_1, s_1 : Y \to Y_1)$ and
$p_2 = (f_2 : X \to Y_2, s_2 : Y \to Y_2)$ and let
$a : Y_1 \to Y_2$ be a morphism of $\mathcal{C}$ such that
$f_2 = a \circ f_1$ and $s_2 = a \circ s_1$.
First assume that $q = (g : Y \to Z', t : Z \to Z')$.
In this case choose a commutative diagram as the one on the left
$$
\vcenter{
\xymatrix{
Y \ar[d]_{s_2} \ar[r]^g & Z' \ar[d]^u \\
Y_2 \ar[r]^h & Z''
}
}
\quad
\Rightarrow
\quad
\vcenter{
\xymatrix{
Y \ar[d]_{s_1} \ar[r]^g & Z' \ar[d]^u \\
Y_1 \ar[r]^{h \circ a} & Z''
}
}
$$
(with $u \in S$),
which implies the diagram on the right is commutative as well.
Using these diagrams we see that both compositions $q \circ p_1$
and $q \circ p_2$ are the equivalence class of
$(h \circ a \circ f_1 : X \to Z'', u \circ t : Z \to Z'')$.
Thus $q \circ p_1 = q \circ p_2$.
The proof of the other case, in which we have to show
$p_1 \circ q = p_2 \circ q$, is omitted. (It is similar to the
case we did.)

\medskip\noindent
Proof of (3). We have to prove associativity of composition.
Consider a solid diagram
$$
\xymatrix{
& & & Z \ar[d] \\
& & Y \ar[d] \ar[r] & Z' \ar@{..>}[d] \\
& X \ar[d] \ar[r] & Y' \ar@{..>}[d] \ar@{..>}[r] & Z'' \ar@{..>}[d] \\
W \ar[r] & X' \ar@{..>}[r] & Y'' \ar@{..>}[r] & Z'''
}
$$
(whose vertical arrows belong to $S$)
which gives rise to three composable pairs.
Using LMS2 we can choose the dotted arrows making the squares commutative
and such that the vertical arrows are in $S$.
Then it is clear that the composition of the three pairs
is the equivalence class of the pair
$(W \to Z''', Z \to Z''')$ gotten by composing the
horizontal arrows on the bottom row and the vertical arrows
on the right column.

\medskip\noindent
We leave it to the reader to check the identity axioms.
\end{proof}

\end{document}
